\documentclass[../../main.tex]{subfiles}
\begin{document}

{\bf[解]}

$y=a\sin x$と$y=1$は$0\le x \le \dfrac{\pi}{2}$にただ一つの共有点$P(t,1)$を持つ．
グラフの概形は下図である．

\begin{figure}[htb]
  \centering
  \begin{tikzpicture}[scale=2.6, >=stealth, every node/.style={font=\small}]

    % --- 1. 定数・パラメータ定義 (極大値 t = pi/3, a = 2*sqrt(3)/3 ≈ 1.155 付近の形状) ---
    % 視認性向上のため、グラフの比率として a = 1.6, t = 40° (rad ≈ 0.698) を採用
    \def\aVal{1.6}
    \def\tVal{0.675}           % a * sin(t) = 1 となる交点の x 座標
    \def\piHalf{1.5708}        % pi / 2

    \coordinate (O) at (0, 0);
    \coordinate (T) at (\tVal, 0);
    \coordinate (P) at (\tVal, 1.0);
    \coordinate (H) at (\piHalf, 0);
    \coordinate (A) at (\piHalf, \aVal);

    % --- 2. 領域のハッチング塗り分け ---
    % (i) S1: y = 1 より上、y = a*sin(x) より下の領域 (水平ストライプ)
    \fill[pattern=horizontal lines, pattern color=blue!70, opacity=0.8]
    (\tVal, 1.0) -- plot[domain=\tVal:\piHalf, samples=50] (\x, {\aVal*sin(\x r)})
    -- (\piHalf, 1.0) -- cycle;

    % (ii) S2: 曲線 y = a*sin(x) (0 <= x <= t) および 直線 y = 1 (t <= x <= pi/2) の下側領域 (斜線)
    \fill[pattern=north east lines, pattern color=teal!70, opacity=0.7]
    (0, 0) -- plot[domain=0:\tVal, samples=40] (\x, {\aVal*sin(\x r)})
    -- (\piHalf, 1.0) -- (\piHalf, 0) -- cycle;

    % --- 3. 座標軸 ---
    \draw[->, thick] (-0.2, 0) -- (2.05, 0) node[right] {$x$};
    \draw[->, thick] (0, -0.2) -- (0, 1.95) node[above] {$y$};
    \node[below left=2pt] at (O) {$O$};

    % --- 4. 曲線 y = a*sin(x) および 直線群 ---
    % 曲線 y = a sin x (0 <= x <= pi/2)
    \draw[very thick, blue!85!black, domain=0:\piHalf, samples=80]
    plot (\x, {\aVal*sin(\x r)});
    \node[blue!85!black, font=\footnotesize, above] at (1.05, 1.6) {$y = a\sin x$};

    % 直線 y = 1
    \draw[thick, black] (-0.1, 1.0) -- (1.9, 1.0);

    % 境界線 x = pi/2
    \draw[dashed, gray!70] (\piHalf, 0) -- (\piHalf, \aVal);

    % --- 5. 補助線と目盛り ---
    % 交点 P(t, 1) からの垂線
    \draw[dashed, gray!70] (\tVal, 0) -- (P);
    \fill[red!85!black] (P) circle (0.7pt) node[above left=1pt, font=\scriptsize] {$P(t,1)$};

    % y 軸上の目盛り (y = 1, y = a)
    \draw[dashed, gray!60] (0, \aVal) -- (\piHalf, \aVal);
    \draw[gray!80!black] (-0.03, 1.0) -- (0.03, 1.0);
    \node[left=2pt] at (0, 1.0) {$1$};
    \draw[gray!80!black] (-0.03, \aVal) -- (0.03, \aVal);
    \node[left=2pt] at (0, \aVal) {$a$};

    % x 軸上の目盛り (x = t, x = pi/2)
    \fill (T) circle (0.5pt) node[below=2pt] {$t$};
    \fill (H) circle (0.5pt) node[below=2pt] {$\frac{\pi}{2}$};

    % --- 6. 領域ラベル S1, S2 ---
    \node[draw=blue!85!black, rounded corners=2pt, inner sep=2pt, font=\footnotesize\bfseries, blue!90!black]
    at (1.18, 1.25) {$S_1$};
    \node[draw=teal!85!black, rounded corners=2pt, inner sep=2pt, font=\footnotesize\bfseries, teal!90!black]
    at (1.05, 0.45) {$S_2$};

  \end{tikzpicture}
  \caption{曲線 $y=a\sin x$ と直線 $y=1$ で分割される2つの領域 $S_1, S_2$}
  \label{fig:sine_curve_regions}
\end{figure}

$P(t,1)$は$y=a\sin t$の上にあるから
\begin{align}
  a\sin t=1\quad \left(0<t\le\dfrac{\pi}{2}\right) \label{eq:1}
\end{align}
を満たす．

$y=a\sin x,y=0,x=\dfrac{\pi}{2}$で囲まれた面積を$S$とすると$S_2 = S - S_1$だから，
題意の面積の差$f=S_2-S_1$は
\begin{align}
  f
  & = S_2 - S_1
  & = (S-S_1)-S_1 \\
  & = S-2S_1 \\
  & = \int_0^{\pi/2}a\sin x\,dx - 2\int_{t}^{\pi/2}\left(a\sin x - 1\right)\,dx \\
  & = -a \left[\cos x \right]_0^{\pi/2} - 2\left[-a\cos x - x\right]_{t}^{\pi/2} \\
  & = a - 2\left(a\cos t - \dfrac{\pi}{2} + t \right) \\
  & = a - 2t - 2a\cos t + \pi \label{eq:2}
\end{align}
である．

これを$t$の関数と見て増減を調べよう．$t$での一階微分を計算すると
\begin{align}
  f'(t)
  &= \dfrac{da}{dt} - 2 - 2\dfrac{da}{dt}\cos t + 2a\sin t \label{eq:3}
\end{align}
だが，\cref{eq:1}の両辺を$t$で微分すると
\begin{align}
  \dfrac{da}{dt} = \dfrac{-\cos t}{\sin^2 t} \label{eq:4}
\end{align}
だから，\cref{eq:1}と合わせて\cref{eq:3}に代入して
\begin{align}
  f'(t)
  &= \dfrac{-\cos t}{\sin^2 t} - 2 + \dfrac{2\cos^2 t}{\sin^2 t} + 2 \\
  &= \dfrac{\cos t\left(2\cos t - 1\right)}{\sin^2 t}
\end{align}
だから，$f(t)$の増減表は以下のようになる．

\begin{table}[htb]
  \centering
  \caption{$f(t)$ の増減表}
  \label{tab:g_increase_decrease}
  \renewcommand{\arraystretch}{1.4}
  \begin{tabular}{c|c|c|c|c|c}
    $t$ & $(0)$ & $\dots$ & $\dfrac{\pi}{3}$ & $\dots$ & $\dfrac{\pi}{2}$ \\ \hline
    $g'(t)$ & & $+$ & $0$ & $-$ & $0$ \\ \hline
    $g(t)$ & & $\nearrow$ & 極大 & $\searrow$ &
  \end{tabular}
\end{table}

従って，求める$f(t)$の最大値は$t=\pi/3$の時の
\begin{align}
  \max f = f\left(\dfrac{\pi}{3}\right) = \dfrac{\pi}{3}
\end{align}
であり，この時\cref{eq:1}から$a$は
\begin{align}
  a=\dfrac{1}{\sin \dfrac{\pi}{3}} = \dfrac{2\sqrt{3}}{3}
\end{align}
である．$\cdots$(答)
\newpage
\end{document}
